\newcommand{\ipPlaceHolder}{
    \foreach \i in {1,...,3}{\TextField[name=ip\i, width=1.5cm,print,borderwidth=0, backgroundcolor={0.9 0.9 0.9}]{}.}\TextField[name=ip4, width=1.5cm,print,borderwidth=0, backgroundcolor={0.9 0.9 0.9}]{}
}

\chapter{Chat-Client mit Python: Schritt-für-Schritt-Anleitung}

In diesem Arbeitsblatt lernen Sie Schritt für Schritt, wie Sie mit Python einen einfachen Chat-Client programmieren, um mit MitschülerInnen Nachrichten über das Internet auszutauschen.

\section{Überblick: Wie funktioniert ein Chat-Client?}

Ein Chat-Client verbindet sich über das Internet mit einem Server, also einem Computer, der Nachrichten von Clients hin- und herschickt. Nachrichten werden als Text über ein das \gls{TCP}-Protokoll verschickt. Der Server kann die Nachricht an andere Clients weiterleiten oder – wie beim Echo-Server – einfach zurückschicken.

\begin{figure}[H]
    \begin{tikzpicture}[
            base/.style={rectangle, draw, minimum width=2cm, minimum height=1cm, text=white, rounded corners, text width=2.4cm, align=center,inner sep=8pt},
            server/.style={base, fill=RoyalBlue},
            client/.style={base, fill=ForestGreen},
            arrow/.style={-Stealth, thick, rounded corners, text width=4cm, align=center}
        ]

        % Nodes
        \node[client] (client1) {Client 1\\\texttt{192.168.0.20}};
        \node[server, right=6cm of client1] (server) {Echo-Server\\\texttt{192.168.0.30}};

        % Arrows
        \draw[arrow] (client1) -- node[above, yshift=4pt] {\texttt{192.168.0.20:Hallo Max!}} (server);
        \draw[arrow] (server) -- node[below, yshift=-4pt] {\texttt{192.168.0.20:Hallo Max!}} (client1);

    \end{tikzpicture}
    \centering
    \caption{Ablauf einer Chat-Nachricht über den Echo-Server}
\end{figure}

\begin{figure}[H]
    \begin{tikzpicture}[
            base/.style={rectangle, draw, minimum width=2cm, minimum height=1cm, text=white, rounded corners, text width=2.4cm, align=center,inner sep=8pt},
            server/.style={base, fill=RoyalBlue},
            client/.style={base, fill=ForestGreen},
            arrow/.style={-Stealth, thick, rounded corners, text width=4cm, align=center}
            ]

            % Nodes
            \node[client] (client1) {Client 1\\\texttt{192.168.0.20}};
            \node[server, below right=of client1] (relay) {Relay-Server\\\texttt{192.168.0.40}};
            \node[client, above right=of relay] (client2) {Client 2\\\texttt{192.168.0.21}};

            % Arrows
            \draw[arrow] (client1) -- node[midway, left] {\texttt{192.168.0.21:Hallo!}} (relay);
            \draw[arrow] (relay) -- node[midway, right] {\texttt{192.168.0.21:Hallo!}} (client2);

    \end{tikzpicture}
    \centering
    \caption{Kommunikation über einen Relay-Server: Die Nachricht wird vom Client an den Relay-Server geschickt und von dort an den Ziel-Client weitergeleitet.}
\end{figure}
Im folgenden werden wir die Kommunikation über einen Relay-Server umsetzen. So können Sie Nachrichten an MitschülerInnen schicken, ohne deren \gls{IP}-Adresse zu kennen.

\section{Aufgaben: Schritt für Schritt zum Chat-Client}

\begin{myexercise}[Meine Öffentliche \gls{IP}-Adresse herausfinden]
Finden Sie die öffentliche \gls{IP}-Adresse Ihres Laptops heraus. Öffnen Sie dazu die Webseite \url{https://whatismyipaddress.com/de/meine-ip} und notieren Sie die \gls{IP}v4-Adresse.

    
Meine öffentliche IP-Adresse ist: \ipPlaceHolder

\end{myexercise}

\begin{myexercise}[IP-Adresse eines Mitschülers/einer Mitschülerin herausfinden]
Fragen Sie einen Mitschüler oder eine Mitschülerin nach seiner/ihrer öffentlichen IP-Adresse und notieren Sie diese.
\begin{myanswer}
Die öffentliche IP-Adresse von \rule{2.5cm}{0.15mm} ist: \ipPlaceHolder
\end{myanswer}
\end{myexercise}

\begin{myexercise}[Nachricht zusammensetzen]
Schreiben Sie ein Python-Programm, das folgende Variablen definiert:
\begin{itemize}
  \item \texttt{ip}: die IP-Adresse Ihres Mitschülers/Ihrer Mitschülerin als Text
  \item \texttt{message}: die Nachricht, die Sie schicken möchten
  \item \texttt{combined}: die kombinierte Chat-Nachricht im Format \texttt{ip:message}
\end{itemize}
Geben Sie den Wert von \texttt{combined} mit \texttt{print()} aus.
\begin{myanswer}
\begin{lstlisting}
ip = "20.161.75.213"
message = "Hallo Bob!"
combined = ip + ":" + message
print(combined)
\end{lstlisting}
\end{myanswer}
\end{myexercise}

\begin{myexercise}[Benutzereingabe und Schleife]

Erweitern Sie Ihr Programm so, dass die IP-Adresse und die Nachricht jeweils per \texttt{input()} eingegeben werden. Kombinieren Sie beide wie oben und geben Sie das Ergebnis aus. Bauen Sie eine Endlosschleife ein, damit Sie beliebig viele Nachrichten verschicken können.
\begin{myanswer}
\begin{lstlisting}[language=Python]
while True:
    ip = input("Empfänger-IP: ")
    message = input("Nachricht: ")
    combined = ip + ":" + message
    print(combined)
\end{lstlisting}
\end{myanswer}
\end{myexercise}

\begin{myexercise}[Server-IP und Port definieren]
Definieren Sie die IP-Adresse und den Port des Chat-Servers als Variablen \lstinline|server_ip| und \lstinline|port| (siehe Tabelle unten).
\begin{myanswer}
\begin{lstlisting}
server_ip = "88.198.193.239"
port = 43223
\end{lstlisting}
\end{myanswer}
\end{myexercise}

\begin{myexercise}[Verbindung zum Server herstellen]

Das Paket \texttt{socket} ist in Python standardmässig bereits installiert und dient dazu, um Netzwerkverbindungen herzustellen. Importieren Sie das Paket und erstellen Sie eine Funktion \lstinline|sende|, die sich mit dem Server verbindet und eine Nachricht sendet. Die Funktion soll die IP-Adresse des Servers, den Port und die Nachricht als Parameter entgegennehmen. Die Antwort des Servers soll ausgegeben werden.

Verwenden Sie folgende Code-Vorlage: Innerhalb Ihrer Definition:
\begin{lstlisting}[language=Python]
import socket

# IHRE DEFINITION HIER
     with socket.socket(socket.AF_INET, socket.SOCK_STREAM) as s: # Socket erstellen und verbinden
        s.connect((server_ip, port)) # Mit Server verbinden
        s.sendall(message.encode()) # Nachricht senden
        data = s.recv(1024) # Antwort empfangen und ausgeben
        print('Antwort vom Server:', data.decode())
\end{lstlisting}

\begin{myanswer}
\begin{lstlisting}[language=Python]
import socket

def send_message(server_ip, port, message):
    with socket.socket(socket.AF_INET, socket.SOCK_STREAM) as s:
        s.connect((server_ip, port))
        s.sendall(message.encode())
        data = s.recv(1024)
        print('Antwort vom Server:', data.decode())
\end{lstlisting}
\end{myanswer}
\end{myexercise}

\begin{myexercise}[Chat-Client zusammenbauen und testen]

Kombinieren Sie alles zu einem vollständigen Chat-Client, so dass Sie in einer Endlosschleife Nachrichten eingeben und an den Server schicken können. Die Antwort des Servers wird ausgegeben.
\begin{myanswer}
\begin{lstlisting}[language=Python]
import socket

server_ip = "116.203.99.55"
port = 43223

while True:
    ip = input("Empfänger-IP: ")
    message = input("Nachricht: ")
    combined = ip + ":" + message
    with socket.socket(socket.AF_INET, socket.SOCK_STREAM) as s:
        s.connect((server_ip, port))
        s.sendall(combined.encode())
        data = s.recv(1024)
        print('Antwort vom Server:', data.decode())
\end{lstlisting}
\end{myanswer}
\end{myexercise}

\begin{figure}[H]
    \begin{tikzpicture}[
            base/.style={rectangle, draw, minimum width=2.2cm, minimum height=1.1cm, text=white, rounded corners, text width=2.5cm, align=center,inner sep=8pt},
            server/.style={base, fill=RoyalBlue},
            client/.style={base, fill=ForestGreen},
            arrow/.style={-Stealth, thick, rounded corners, text width=4.2cm, align=center}
        ]

        % Nodes
        \node[client] (client) {Client\\\texttt{192.168.0.20}};
        \node[server, right=9cm of client] (server) {Echo-Server\\\texttt{192.168.0.30}};

        % Arrows
        \draw[arrow] (client) -- node[above, yshift=4pt] {\texttt{192.168.0.20:Hallo Max!}} (server);
        \draw[arrow] (server) -- node[below, yshift=-4pt] {\texttt{192.168.0.20:Hallo Max!}} (client);

    \end{tikzpicture}
    \centering
    \caption{Ablauf einer Chat-Nachricht über den Echo-Server}
\end{figure}

\noindent
Der Client schickt eine Nachricht im Format \texttt{IP:Nachricht} an den Server. Der Echo-Server schickt exakt dieselbe Nachricht zurück. So können Sie testen, ob Ihr Client korrekt funktioniert.
